\documentclass[english, parskip=full]{scrartcl}
\usepackage[utf8]{inputenc}  % Kodierung der Datei
\usepackage[english]{babel}  % Sprache auf Deutsch 
\usepackage{color}
\usepackage{graphicx, gensymb}
\usepackage{amsfonts, cancel, amssymb}
\usepackage{amsmath, amsthm, array, esint}
\usepackage{booktabs, multirow}  % schönere Tabellen
\usepackage{caption}  % Anpassung der Captions
\usepackage{scrlayer-scrpage}    % Manipulation des Seitenstils
\pagestyle{scrheadings}  % Stil für die Seite setzen
\automark{section}  % Abschnittsnamen als Seitenbeschriftung 
\setheadsepline{.5pt}    % Kopzeile durch Linie abtrennen
\usepackage[hidelinks]{hyperref}  % Links und weitere PDF-Features
\usepackage{typearea, tikz, pgffor, verbatim}
\usetikzlibrary{calc, quotes}
\usepackage[cache=false, outputdir=build]{minted}
\usepackage{placeins} % provides \FloatBarrier

\definecolor{bg}{rgb}{0.9,0.9,0.9}
\newcommand\numberthis{\addtocounter{equation}{1}\tag{\theequation}}
\newcommand{\bra}{}
\newcommand{\kom}{}
\newcommand{\brak}{}
\newcommand{\braket}{}
\newcommand{\intx}{}
\newcommand{\intr}{}
\newcommand{\kla}{}
\newcommand{\klam}{}
\newcommand{\klammer}{}
\newcommand{\oper}{}
\newcommand{\tecxt}{}
\newcommand{\Bra}{}
\newcommand{\Ket}{}
\renewcommand{\i}{\text{i}}
\newcommand{\e}{\text{e}}
\newcommand*{\Fourier}{\textsc{Fourier}\ }
\newcommand*{\F}{\mathcal{F}}
\newcommand*{\sinc}{\text{sinc\,}}
\newcommand{\conv}{\mathop{\scalebox{3.5}{\raisebox{-0.38ex}{\makebox[0.5\width][c]{$\ast$}}}}\limits}

\newcommand*{\titel}{04 Partition functions}
\newcommand*{\autor}{Joachim Schwardt}

\hypersetup{pdfauthor={\autor}, pdftitle={\titel}}

%\titlehead{titlehead}
\subject{Theoretical Physics (Master)}
\title{\titel}
\author{\autor}
\date{\today}

\begin{document}

%\begin{titlepage}
\maketitle  % Titel setzen
%\tableofcontents  % Inhaltsverzeichnis setzen
%\end{titlepage}

\section*{Problem 1: \normalfont\textit{\large{Partition functions}}}
We calculate the partition function of the harmonic oscillator in two different ways.
%\subsection*{a)}
The direct definition is
\begin{align}
Z &:= \sum_{n\ge 0}\e^{-\beta E_n} = \sum_{n\ge 0}\e^{-\beta\omega \kla\left( n+\frac{1}{2} \right)} = \e^{-\beta\omega/2} \frac{1}{1 - \e^{-\beta\omega}} = \frac{2}{\sinh\kla\left( \frac{\beta\omega}{2} \right)}.
\end{align}
%\subsection*{b)}
The propagator for the harmonic oscillator is
\begin{align}
G &= \sqrt{\frac{m\omega}{2\pi\i\sin (\omega t)}} \exp\kla\left( \frac{\i m \omega}{2\sin (\omega t)}\klam\left[ (x_N^2 + x_0^2)\cos(\omega t) - 2x_Nx_0 \right] \right).
\end{align}
The partition function may also be computed from
\begin{align}
Z &= \intr\int_{\mathbb{R}^{ }} G(x, t=-\i\beta | x, 0) \, \mathrm{d}^{ }x = \\
&= \sqrt{\frac{m\omega}{2\pi\i\sin(-\omega \i\beta)}} \intr\int_{\mathbb{R}^{ }} \exp\kla\left( \frac{\i m\omega}{\sin (-\omega \i\beta)} x^2\klam\left[ \cos(\omega \i\beta) - 1 \right] \right) \, \mathrm{d}^{ }x = \\
&= \sqrt{\frac{m\omega}{2\pi\i\sin(-\omega \i\beta)}} \sqrt{\frac{\pi \sin(-\omega\i\beta)}{-\i m\omega \kla\left( \cosh(\omega\beta) - 1 \right)}} = \\
&= \frac{2}{\sqrt{\frac{\cosh(\beta\omega) - 1}{2}}} = \frac{2}{\sinh\kla\left( \frac{\beta\omega}{2} \right)}.
\end{align}
Here we used
\begin{align}
\sinh^2(x) = \frac{\e^{2x} + \e^{-2x} - 2}{4} = \frac{\cosh(2x)}{2} - \frac{1}{2}
\end{align}
with $x = \frac{\beta\omega}{2}$.

\newpage
\section*{Problem 2: \normalfont\textit{\large{Single-mode Bogoliubov transforms}}}
Consider a quadratic bosonic Hamiltonian of the form (this could be a bosonic toy model for the Bardeen-Cooper-Schrieffer theory of superconductivity)
\begin{align}
H &= \omega a^\dagger a + \eta a^2 + \eta^\ast a^{\dagger\,2}.
\end{align}
\subsection*{a)}
Let $a = ub + vb^\dagger$, where we want $b$ to be a bosonic annihilation operator.
Then we find
\begin{align}
\klam\left[ a,a^\dagger  \right] &= \klam\left[ ub + vb^\dagger, u^\ast b^\dagger + v^\ast b \right] = uu^\ast - vv^\ast \overset{!}{=} 1,\\
\klam\left[ a,a \right] &= \klam\left[ ub + vb^\dagger, ub + vb^\dagger \right] = uv - vu = 0.
\end{align}

\subsection*{b)}
Now let $u,v\in\mathbb{R}$.
Then we may express the Bogoliubov as 
\begin{align}
\begin{pmatrix}
a \\ a^\dagger
\end{pmatrix} &= \begin{pmatrix}
u & v \\ v & u
\end{pmatrix}\begin{pmatrix}
b \\ b^\dagger
\end{pmatrix}.
\end{align}

\subsection*{c)}


\section*{Problem 3: \normalfont\textit{\large{Coherent state of Landau levels}}}
The Hamiltonian of an electron sensing a magnetic field is
\begin{align}
H &= \frac{p_x^2 + p_y^2}{2m} + \frac{m}{2}\omega_L^2 (x^2+y^2) - \omega_L (xp_y-yp_x) = \\
&= \omega_L\kla\left( a_x^\dagger a_x + a_y^\dagger a_y + 1 \right) + \i\omega_L\kla\left( a_x^\dagger a_y - a_y^\dagger a_x \right) = 2\omega_L\kla\left( b^\dagger b + \frac{1}{2} \right)
\end{align}
where $\omega_L = \frac{eB}{2m}$ is the Larmor frequency.
Now consider a state of the form
\begin{align}
\Ket | \Psi \rangle &:= D_+(\alpha_+)D_-(\alpha_-) \Ket | 0 \rangle = \exp\kla\left( \alpha_+b_+^\dagger - \alpha_+^\ast b_+ \right)  \exp\kla\left( \alpha_-b_-^\dagger - \alpha_-^\ast b_- \right) \Ket | 0 \rangle.
\end{align}
\subsection*{a)}
Then we may compute the expectation values
\begin{align}
\bra\langle x \rangle &= 
\end{align}
\subsection*{b)}

\section*{Problem 4: \normalfont\textit{\large{Boundary effects on Landau levels}}}
Consider the Landau Hamiltonian with an additional effective boundary at $x=0$, i.e. the potential
\begin{align}
V(x) &= \left\{\begin{matrix}
\infty & : & x<0 \\ 0 & : &  x\ge 0
\end{matrix} \right. .
\end{align}
In the Landau gauge $\vec{A} = Bx\vec{e}_y$, the full Hamiltonian then becomes
\begin{align}
H &= \frac{1}{2m}\klam\left[ p_x^2 + \kla\left( p_y - eBx \right)^2 \right] + V(x).
\end{align}
\subsection*{a)}
We can write the Hamiltonian in the more explicit form
\begin{align}
H &= \frac{p_x^2+p_y^2}{2m} + \frac{e^2B^2x^2}{2m} - \frac{eB}{2m}\klammer\left\{ p_y, x \right\} + V(x) = \\
&=  \frac{p_x^2+p_y^2}{2m} + \frac{m}{2}\omega_B^2x^2 - \omega_Bp_yx + V(x).
\end{align}
With $k_y$ being the eigenvalue of $p_y$, we may express our Hamiltonian with the effective potential
\begin{align}
V_\tecxt\text{eff}(x) &= \frac{m}{2}\omega_B^2x^2-\omega_Bk_yx + V(x).
\end{align}
\subsection*{b)}
For $k_y=0$ the Hamiltonian effectively becomes a one dimensional harmonic oscillator potential, but with the additional requirement $\psi(x=0)=0$.
Thus only the odd eigenfunctions $\psi_k(x)$ for $k=2n+1$ are possible, which leads to the energies
\begin{align}
E_n(k_y=0) &= \omega_b \kla\left( [2n+1] + \frac{1}{2} \right).
\end{align}
For $k_y\rightarrow\infty$ ???
\subsection*{c)}

\section*{Problem 5: \normalfont\textit{\large{Tunneling probability}}}
Consider a particle subject to the potential
\begin{align}
V(x) = \alpha (x^2 - \Delta^2)^2
\end{align}
with $\alpha,\Delta> 0$.
\subsection*{a)}
\subsection*{b)}
Now consider the negative of the potential $V$.
Then conservation of energy gives
\begin{align}
E &= \frac{m}{2}\dot{x}^2 + \alpha (x^2-\Delta^2)^2.
\end{align}
Separation of variables leads to
\begin{align}
\pm\sqrt{\frac{2}{m}} t &= \intx\int_{0}^{t}\frac{\dot{x}}{\sqrt{E - \alpha (x^2 - \Delta^2)^2}}\, \mathrm{d}t = \intx\int_{x_0}^{x}\frac{1}{\sqrt{E - \alpha (s^2 - \Delta^2)^2}}\, \mathrm{d}s \kom\overset{\substack{{s\,=\,\Delta\cosh (u)} \\ |}}{=} \\
&= \intx\int_{ }^{ }\frac{\Delta \sinh (u)}{\sqrt{E - \alpha\Delta^4 \sinh^4(u)}}\, \mathrm{d}u = \sqrt{\frac{\alpha\Delta^4}{E}}
\end{align}
Now substitute $u := \kla\left( E - \alpha(s^2-\Delta^2)^2 \right)^{-1/2}$.
Then we find $s = \sqrt{\Delta^2 + \sqrt{\frac{1}{\alpha}\kla\left( E - \frac{1}{u^2} \right)}}$ and the differential relation
\begin{align}
\frac{\tecxt\text{d}u}{\tecxt\text{d}s} &= -\frac{1}{2}u^3\kla\left( -2\alpha (s^2-\Delta^2)(2s) \right) = \\
&= 2\alpha u^3 \sqrt{\Delta^2 + \sqrt{\frac{1}{\alpha}\kla\left( E - \frac{1}{u^2} \right)}}\sqrt{\frac{1}{\alpha}\kla\left( E - \frac{1}{u^2} \right)}
\end{align}
The differential equation is
\begin{align}
m\ddot{x} &= -V'(x) = 4\alpha x (x^2-\Delta^2) &&\Rightarrow & \ddot{x} + \omega^2x &= \frac{4\alpha}{m} x^3
\end{align}
with $\omega^2 := \frac{4\alpha\Delta^2}{m} $ and .
Now measure time with the dimensionless parameter $s := \omega t$ to obtain $\ddot{x}(t) = \ddot{x}(s(t)) = x''(s) \dot{s}^2 = x''(s)\omega^2$ and
\begin{align}
x'' + x &= ax^3
\end{align}
where $a := \frac{4\alpha}{m\omega^2}$.
This is the anharmonic oscillator, and it is not solvable analytically.
However, we may very well be able to obtain the solution as an asymptotic series.
Start with $a=0$, which yields the unperturbed
\begin{align}
x_0(s) &= x_0\cos (rs).
\end{align}
Here we chose the cosine solution, since we know that the initial and final position are supposedly antisymmetric with $x(t_\tecxt\text{in}=0) = x_\tecxt\text{in} = -\Delta$ and $x(t_\tecxt\text{out}=T) = x_\tecxt\text{out} = \Delta$.
Then we can write the full solution as a series in powers of $a$, i.e.
\begin{align}
x(s) &= \sum_{k\ge 0} x_k(s)a^k, 
\end{align}
where we have already determined the zeroth coefficient.
We may insert the initial conditions at this point, giving
\begin{align}
x_0(s) = -\Delta \cos (s)
\end{align}
where $\cos(S := \omega T) \overset{!}{=} -1$ gives $T = \frac{\pi}{\omega}$.
Note that
\begin{align}
x^3 &= \kla\left( \sum_k x_ka^k \right)^3 = \sum_ka^k\sum_{i+j+l=k} x_ix_jx_l
\end{align}
Inserting our series into the differential equation yields
\begin{align}
\sum_ka^k\kla\left( x_k''+x_k \right) &= \sum_ka^{k+1}\sum_{i+j+l=k} x_ix_jx_l.
\end{align}
For $k=1$ we find :
\begin{align}
x_1''+x_1 &= 3x_0^2x_1 = 3\Delta^2\cos^2(s)x_1
\end{align}

%\begin{thebibliography}{9}
%\bibitem{example} description, \url{https://www.exmapleurl}, day.\,month~2021.
%\end{thebibliography}
\end{document}